% GENERUOJAMA is rezultatai/darbiniai/sumaisymas_*.csv
% Ranka NELIESTI - paleisti: python -m src.eksperimentai.klaidos
\begingroup
\footnotesize
\setlength{\tabcolsep}{5pt}
\begin{tabularx}{\textwidth}{@{}>{\raggedright\arraybackslash}X>{\raggedleft\arraybackslash}p{1.6cm}>{\centering\arraybackslash}p{1.9cm}>{\centering\arraybackslash}p{1.9cm}>{\centering\arraybackslash}p{1.9cm}@{}}
\toprule
\textbf{Kategorija} & \textbf{$n$} & \textbf{Tikslumas} & \textbf{Atkūrimas} & \textbf{F1} \\
\midrule
\multicolumn{5}{@{}l}{\textbf{XGBoost}} \\
\quad DDoS & 157\,500 & 0,955 & 0,787 & 0,863 \\
\quad Recon & 55\,691 & 0,895 & 0,772 & 0,829 \\
\quad DoS & 55\,266 & 0,596 & 0,896 & 0,716 \\
\quad Mirai & 45\,000 & 0,998 & 0,998 & 0,998 \\
\quad Spoofing & 30\,000 & 0,925 & 0,878 & 0,901 \\
\quad Benign & 15\,000 & 0,543 & 0,787 & 0,643 \\
\quad Web & 3\,556 & 0,312 & 0,486 & 0,380 \\
\quad BruteForce & 1\,878 & 0,396 & 0,492 & 0,439 \\
\addlinespace
\multicolumn{5}{@{}l}{\textbf{Random Forest}} \\
\quad DDoS & 157\,500 & 0,954 & 0,787 & 0,862 \\
\quad Recon & 55\,691 & 0,859 & 0,829 & 0,844 \\
\quad DoS & 55\,266 & 0,595 & 0,893 & 0,714 \\
\quad Mirai & 45\,000 & 0,998 & 0,998 & 0,998 \\
\quad Spoofing & 30\,000 & 0,942 & 0,868 & 0,903 \\
\quad Benign & 15\,000 & 0,552 & 0,748 & 0,635 \\
\quad Web & 3\,556 & 0,423 & 0,362 & 0,390 \\
\quad BruteForce & 1\,878 & 0,559 & 0,411 & 0,474 \\
\addlinespace
\multicolumn{5}{@{}l}{\textbf{MLP}} \\
\quad DDoS & 157\,500 & 0,953 & 0,762 & 0,847 \\
\quad Recon & 55\,691 & 0,936 & 0,558 & 0,699 \\
\quad DoS & 55\,266 & 0,571 & 0,892 & 0,696 \\
\quad Mirai & 45\,000 & 0,995 & 0,998 & 0,996 \\
\quad Spoofing & 30\,000 & 0,911 & 0,778 & 0,839 \\
\quad Benign & 15\,000 & 0,483 & 0,691 & 0,568 \\
\quad Web & 3\,556 & 0,125 & 0,617 & 0,208 \\
\quad BruteForce & 1\,878 & 0,137 & 0,604 & 0,223 \\
\bottomrule
\end{tabularx}
\vspace{2pt}
\raggedright\scriptsize Rodikliai išmatuoti \textbf{testavimo} aibėje ties \emph{argmax}, o ne ties klaidingų teigiamų biudžetu, todėl su pagrindine aptikimo kokybės lentele tiesiogiai negretinami. $n$ --- tikrų pavyzdžių skaičius; retose kategorijose vienas kitoks sprendimas keičia atkūrimą per dešimtąsias procentinio punkto dalis.
\endgroup
